\begin{table}[!htbp]
\centering
\footnotesize
\makebox[\textwidth][c]{%
\begin{threeparttable}
\caption{US House district-level RMSE by sample size and race competitiveness.}
\label{tab:us_house_rmse_grid}
\begin{tabular}{lccccc}
\toprule
& \multicolumn{5}{c}{Winner's vote share (\%)} \\
\cmidrule(lr){2-6}
Min. respondents & 20--40 & 40--60 & 60--80 & 80--90 & 90--100 \\
\midrule
\multicolumn{6}{l}{RMSE (pp)} \\
\addlinespace[0.2em]
All & 10.0 & 8.7 & 9.3 & 10.5 & 10.4 \\
30+ & 10.0 & 8.7 & 9.2 & 10.6 & 9.8 \\
60+ & 4.3 & 8.2 & 8.6 & 9.8 & 10.5 \\
90+ & -- & 7.6 & 7.4 & 10.0 & 12.4 \\
120+ & -- & 6.9 & 6.8 & 9.2 & 16.6 \\
\midrule
\multicolumn{6}{l}{Number of district-years} \\
\addlinespace[0.2em]
All & 5 & 1256 & 1799 & 195 & 156 \\
30+ & 5 & 1251 & 1764 & 186 & 148 \\
60+ & 3 & 1060 & 1434 & 129 & 89 \\
90+ & 0 & 642 & 681 & 48 & 28 \\
120+ & 0 & 219 & 199 & 9 & 3 \\
\bottomrule
\end{tabular}
\begin{tablenotes}
\footnotesize
\item \textit{Note:} Each cell reports RMSE (in percentage points) for all congressional district-years (2006--2022) meeting both the respondent threshold (row) and winner's vote share range (column). A district-year is one district in one election year, so the same district recurs across years. RMSE is pooled over all qualifying observations rather than averaged across years, so years with more observations count for more. Estimates use Matching + CES Weighting.
\end{tablenotes}
\end{threeparttable}%
}
\end{table}